\makenoidxglossaries

\newglossaryentry{downscaling}{
    name=downscaling,
    description={The method of reducing distance between measurements in climate models by increasing the amount of data points}
}

\newglossaryentry{realisation}{
    name=realisation,
    description={GCMs and RCMs are typically run multiple times creating ensembles. One sample of a such ensemble is called a realisation}
}

\newacronym
[description={Global climate models, that estimate the climate on the entire planet on a long temporal scale. Often in a 0.5° to 1.25\textdegree resolution}]
{gcm}{GCM}{General Circulation Model}


\newacronym
[description={Model created for downscaling the output General Circulation Models to a regional or local domain, usually improving specificity}]
{rcm}{RCM}{Regional Climate Model}

\newacronym
[description={Methods for \gls{downscaling}, in which data provides the general understanding (as opposed to physics). These methods are often Machine Learning based}]
{sdm}{SDM}{Statistical Downscaling Method}

\newacronym
[description={This is downscaling model similar to an SDM, but specifically this term covers physics-based downscaling models as opposed to the statistics-based SDM}]
{ddm}{DDM}{Dynamical Downscaling Model}

\newacronym
[description={The best of the best. Crème de la crème. The, until now, best method or model for a specific task}]
{sota}{SOTA}{State Of The Art}